\appendix

\section{Threshold Sensitivity}\label{app:sensitivity}

The triage thresholds ($\tau_{\text{noise}} = 0.12$, $\tau_{\text{cat2}} = 0.4$)
and the triage-weighting upweight factor ($w = 2.0$, the sample-weight multiplier
that the triage-weighting strategy of \S\ref{sec:method:menu} applies to the correct
same-class neighbours of Cat\,2 errors) are fixed across all datasets. We swept each threshold across a
range of values and measured downstream XGBoost accuracy: $\tau_{\text{noise}}
\in [0.05, 0.30]$, $\tau_{\text{cat2}} \in [0.2, 0.6]$, $w \in [1.0, 5.0]$.
Across the benchmark datasets, 74\% show $< 0.5\%$ accuracy
variation across all configurations, supporting the use of fixed thresholds
as a sensible default. The two methods (clean-masked SMOTE, triage
weighting) maintain their relative ordering vs.\ baselines at every threshold
configuration tested.

\section{Out-of-Fold Triage Validation}\label{app:oof}

The triage categories underpin the mechanistic argument (\S\ref{sec:results:interventional}),
so we verify they are not optimistic in-sample artefacts. Two facts establish this.
First, by construction every model-derived signal that drives the categorization is
\emph{out-of-bag}: the error indicator is the OOB ensemble majority vote, and the
true-class probability is the OOB estimate (\S\ref{sec:method:uncert}); the only
non-model signal, the local class ratio, is a property of the data geometry.
Second, as a strict check we re-derive the categories under inner $5$-fold
cross-validation: forests fit on the inner-training folds categorize the held-out
instances (via held-out predictions for the error indicator, true-class
probability, and the full aleatoric/epistemic decomposition; neighbour-based
signals are measured against the inner-training set). Over the 79 datasets, the
Cat\,3 (irreducible) fraction among errors is essentially unchanged: the in-sample
and out-of-fold means coincide at $0.838$, the per-dataset Spearman correlation is
$0.97$ ($p<10^{-40}$), and the mean absolute shift is $\OOFCATSHIFT$\,pp (median
$1.0$\,pp). The aleatoric/epistemic ordering that motivates the categories is
preserved out-of-fold (Cat\,3 aleatoric uncertainty exceeds Cat\,2's on $80\%$ of
datasets). The category \emph{identification} is therefore not an in-sample
artefact. We note separately that the small real-data \emph{interventional} effect
(\S\ref{sec:results:interventional}) is sensitive to this re-estimation: with
out-of-fold seeds and a reduced $3\times3$ protocol the targeted-augmentation
contrast remains directionally consistent (random augmentation null) but is no
longer individually significant, which is why we rest the mechanism on the
controlled experiment (\S\ref{sec:results:mechanism}) and the
discrimination/threshold analysis (\S\ref{sec:results:frontier}) rather than on the
real-data intervention.

\section{GBM Triage Agreement}\label{app:gbm}

To verify that the three-way decomposition is not specific to random
forests, we replaced the RF ensemble with bagged XGBoost (5 bags of 100
trees, matching the RF configuration). On a 20-dataset agreement subset,
the GBM-based triage assigns 82\% of training instances to the same category
as the RF-based triage. Downstream accuracy under both triages is
comparable (mean absolute difference $0.04\%$), indicating that any
ensemble model exposing per-tree predictions can drive the decomposition.

\section{Subsampling Category-Structure Control}\label{app:subsampling}

We verified that stratified subsampling to 10\,000 instances preserves the
triage category distribution. On the five largest datasets in our
benchmark, we computed the per-category fraction $\{p_{\text{correct}},
p_{\text{noise}}, p_{\text{data-limited}}, p_{\text{irreducible}}\}$ on the
full dataset and on a single stratified 10\,000-instance subsample. The
mean absolute deviation between the two distributions is $< 2$ percentage
points per category across all five datasets, validating the subsampling
protocol used throughout (\S\ref{sec:results:frontier}).
Detailed per-dataset numbers are provided in the supplementary material.

\section{Computational Cost Detail}\label{app:cost}

Full cost-table breakdown by dataset (extending \S\ref{sec:results:cost}):

\begin{table}[t]
\centering
\caption{Triage wall-clock cost vs.\ a single XGBoost fit (mean over repeats). Overhead = triage time as a percentage of total pipeline time. The triage is a one-time, classifier-agnostic preprocessing cost.}
\label{tab:cost}
\begin{tabular}{lrrrrr}
\toprule
Dataset & $n$ & $d$ & Triage (s) & XGB (s) & Overhead (\%) \\
\midrule
\texttt{iris} & 150 & 4 & 1.04 & 0.15 & 87.2 \\
\texttt{wine} & 178 & 13 & 1.31 & 0.08 & 94.1 \\
\texttt{sonar} & 208 & 60 & 1.10 & 0.12 & 90.2 \\
\texttt{glass} & 214 & 9 & 1.17 & 0.96 & 64.8 \\
\texttt{ionosphere} & 351 & 34 & 1.05 & 0.12 & 89.9 \\
\texttt{vehicle} & 846 & 18 & 1.27 & 0.40 & 76.1 \\
\texttt{vowel} & 990 & 12 & 1.91 & 1.59 & 56.0 \\
\texttt{credit-g} & 1000 & 20 & 1.43 & 0.16 & 89.8 \\
\texttt{segment} & 2310 & 19 & 1.56 & 0.93 & 63.4 \\
\texttt{phoneme} & 5404 & 5 & 2.71 & 0.74 & 77.6 \\
\texttt{satimage} & 6430 & 36 & 2.87 & 2.19 & 58.5 \\
\texttt{eeg-eye-state} & 14980 & 14 & 10.20 & 0.60 & 94.6 \\
\texttt{MagicTelescope} & 19020 & 10 & 8.40 & 1.25 & 88.8 \\
\texttt{bank-marketing} & 45211 & 16 & 16.79 & 0.84 & 95.6 \\
\texttt{covertype} & 50000 & 54 & 32.63 & 49.05 & 50.1 \\
\bottomrule
\end{tabular}
\end{table}


The triage time scales roughly as $O(n \log n)$ in the number of training
instances and is dominated by the random-forest fits; the local class-ratio
computation (k-NN with $k=5$) contributes $< 5\%$ of total triage time
across all dataset sizes.

\section{Instrument Robustness: Non-Tree Ensembles}\label{app:nontree}

The remedy decomposition's model-derived signals come from a random-forest
ensemble (\S\ref{sec:method:uncert}). To verify the headline is not an
artefact of tree-based uncertainty, we re-derive the full decomposition with
two bagged non-tree instruments over the full 27 KEEL + 54
OpenML/scikit-learn roster (81 dataset cells; the headline statistics of
Table~\ref{tab:reducibility} use the 79-dataset menu roster of \S\ref{sec:design}): $B{=}25$
bootstrap-bagged multilayer perceptrons (scaled, $64{\times}32$ hidden units,
early stopping) and $B{=}25$ bootstrap-bagged logistic regressions. Every
signal remains out-of-bag: only ensemble members whose bootstrap excluded an
instance contribute to its probabilities, and the Shaker-style
aleatoric--epistemic decomposition is taken over those members. The local
class ratio is the identical $k$-NN computation, and the error set, one-vs-rest
intervention, and cross-fitted threshold follow \S\ref{sec:method:remedy}
exactly. Two design notes: (i)~members
are unweighted, matching the headline instrument (\S\ref{sec:method:remedy}) ---
class-balanced members enact the operating-point shift
\emph{inside} the instrument (completely so for logistic regression,
\S\ref{sec:results:frontier}) and dissolve the default-threshold deficit being
decomposed; (ii)~Cat\,1 is gated by TCP alone, as the forest-consensus noise
score has no non-tree analogue, which shifts part of the residue from the
data-reducible/irreducible classes into noise but leaves the
threshold-recoverability axis untouched.

\begin{table}[t]
\centering
\caption{Remedy decomposition by instrument (means over datasets with
$\ge 5$ minority errors; definitions of \S\ref{sec:method:remedy}). The
threshold-recoverable share, the paper's headline, is stable across
instruments; the per-dataset correlation with the RF instrument is given in
the last column. The non-tree instruments assign more of the residue to noise
(TCP-only gate, note (ii)).}
\label{tab:nontree}
\footnotesize
\begin{tabular}{lrrrrrr}
\toprule
Instrument & thr.-rec. & (median) & data-red. & irred. & noise & $r$ vs RF \\
\midrule
Random forests (paper instrument) & 69\% & (75\%) & 7\% & 16\% & 7\% & --- \\
Bagged MLPs & 69\% & (70\%) & 3\% & 11\% & 18\% & 0.68 \\
Bagged logistic regression & 68\% & (74\%) & 1\% & 10\% & 21\% & 0.55 \\
\bottomrule
\end{tabular}

\end{table}

Threshold-recoverable remains the dominant remedy class under both non-tree
instruments (Table~\ref{tab:nontree}): $62\%$ for bagged MLPs and $63\%$ for
bagged logistic regression against $68\%$ for the RF instrument, rising to
$73\%$ on $\mathrm{IR}>3$ datasets for both (RF: $80\%$), with per-dataset
correlations of $r=0.78$ and $0.72$ against the forest instrument. Part of the
gap to the RF share is the TCP-only noise gate of note (ii), which absorbs
$19$--$24\%$ of the residue into noise. The decomposition's ordering,
threshold-recoverable far above every other remedy class, is
instrument-independent.

\section{Ensemble-Size Sensitivity}\label{app:msens}

The instrument uses $M{=}5$ forests of $T{=}100$ trees, matching standard
deep-ensemble practice ($M{=}5$ in \cite{lakshminarayanan2017simple}).
Recomputing the full decomposition for $M \in \{1,2,5,10,20\}$ on a
10-dataset IR-spanning subset (Table~\ref{tab:msens}): the triage category
fractions move by at most $2.6$\,pp between $M{=}5$ and $M{=}20$, and the
threshold-recoverable share by $3.3$\,pp on average ($1.9$\,pp for $M{=}10$);
the largest single-dataset deviations occur where the minority error set is
tiny (a handful of instances, so one flip moves the fraction by many points).
Below $M{=}5$ the decomposition is visibly unstable (single-dataset swings up
to $34$\,pp). $M{=}5$ therefore sits at
the stability knee, and the conclusions do not change with a larger ensemble.

\begin{table}[t]
\centering
\caption{Ensemble-size sensitivity: mean threshold-recoverable share and
absolute deviations from the $M{=}5$ reference across 10 datasets.}
\label{tab:msens}
\footnotesize
\begin{tabular}{lrrrr}
\toprule
$M$ & mean thr.-rec. & mean $|\Delta$thr.-rec.$|$ & max $|\Delta$thr.-rec.$|$ & max $|\Delta$cat.\ fraction$|$ \\
\midrule
1 & 66\% & 9.4\,pp & 28.6\,pp & 6.4\,pp \\
2 & 71\% & 7.0\,pp & 42.9\,pp & 7.5\,pp \\
5 (ref.) & 70\% & 0.0\,pp & 0.0\,pp & 0.0\,pp \\
10 & 68\% & 3.7\,pp & 11.2\,pp & 2.3\,pp \\
20 & 69\% & 2.3\,pp & 6.2\,pp & 2.3\,pp \\
\bottomrule
\end{tabular}

\end{table}

\section{Extended Oversampler Comparison}\label{app:oversamplers}

Table~\ref{tab:smote} reports the full oversampler family, including the
Kov\'acs top-ranked variants and Geometric SMOTE (added at review), against
standard SMOTE on the OpenML roster.

\begin{table}[t]
\centering
\caption{Oversampler family vs.\ standard SMOTE (XGBoost, 52 OpenML-roster datasets): mean metric, win rate over SMOTE, and Wilcoxon $p$. The geometric variant G-SMOTE (added at review) buys significantly \emph{less} balanced accuracy than SMOTE at a comparable accuracy level; only polynom-fit (a known weak-tradeoff variant) and the triage clean-masked variant significantly recover accuracy.}
\label{tab:smote}
\begin{tabular}{lrrrrrr}
\toprule
Method (vs.\ SMOTE) & Mean Acc & WR Acc & $p_{\text{Acc}}$ & Mean BAcc & WR BAcc & $p_{\text{BAcc}}$ \\
\midrule
baseline & 0.8538 & 53.8\% & 0.001** & 0.7768 & 11.5\% & 0.000*** \\
borderline\_smote & 0.8505 & 38.5\% & 0.914 & 0.7911 & 25.0\% & 0.013* \\
adasyn & 0.8507 & 26.9\% & 0.180 & 0.7909 & 32.7\% & 0.444 \\
safe\_level\_smote & 0.8484 & 36.5\% & 0.379 & 0.7900 & 34.6\% & 0.282 \\
polynom\_fit\_smote & 0.8542 & 55.8\% & 0.001*** & 0.7800 & 17.3\% & 0.000*** \\
prowsyn & 0.8527 & 48.1\% & 0.128 & 0.7866 & 21.2\% & 0.000*** \\
mwmote & 0.8517 & 38.5\% & 0.771 & 0.7856 & 28.8\% & 0.001** \\
gsmote & 0.8510 & 44.2\% & 0.677 & 0.7814 & 25.0\% & 0.000*** \\
napierala\_guided\_smote & 0.8507 & 51.9\% & 0.103 & 0.7885 & 32.7\% & 0.036* \\
clean\_masked\_smote & 0.8519 & 46.2\% & 0.016* & 0.7863 & 17.3\% & 0.000*** \\
\bottomrule
\end{tabular}
\end{table}


\clearpage
\section{Dataset Listing}\label{app:datasets}

Table~\ref{tab:datasets} lists the OpenML/scikit-learn roster, including
each dataset's OpenML ID, number of instances, features, classes, and imbalance
ratio (majority-class count divided by minority-class count); the remaining 27
datasets are the standard \texttt{imbalanced-learn}/KEEL benchmark suite.

\begin{longtable}{lrrrrc}
\caption{Full dataset listing (54 datasets; the recovered original-paper roster). $n$ = instances; $d$ = features; $C$ = classes; IR = imbalance ratio (majority / minority count). ``sklearn'' marks datasets loaded from scikit-learn rather than OpenML; \texttt{webpage}'s class statistics are not reported because its sparse high-dimensional format excludes it from the menu evaluation (\S\ref{sec:design}).}\label{tab:datasets}\\
\toprule
Dataset & OpenML ID & $n$ & $d$ & $C$ & IR \\
\midrule
\endfirsthead
\multicolumn{6}{l}{\tablename~\thetable\ -- continued}\\
\toprule
Dataset & OpenML ID & $n$ & $d$ & $C$ & IR \\
\midrule
\endhead
\midrule \multicolumn{6}{r}{\textit{continued on next page}}\\
\endfoot
\bottomrule
\endlastfoot
\texttt{phoneme} & 1489 & 5404 & 5 & 2 & 2.4 \\
\texttt{electricity} & 151 & 45312 & 8 & 2 & 1.4 \\
\texttt{bank-marketing} & 1461 & 45211 & 16 & 2 & 7.5 \\
\texttt{MagicTelescope} & 1120 & 19020 & 10 & 2 & 1.8 \\
\texttt{MiniBooNE} & 41150 & 50000 & 50 & 2 & 2.6 \\
\texttt{jannis} & 41168 & 50000 & 54 & 4 & 23.2 \\
\texttt{covertype} & 1596 & 50000 & 54 & 7 & 109.5 \\
\texttt{credit-g} & 31 & 1000 & 20 & 2 & 2.3 \\
\texttt{vehicle} & 54 & 846 & 18 & 4 & 1.1 \\
\texttt{segment} & 36 & 2310 & 19 & 7 & 1.0 \\
\texttt{kc1} & 1067 & 2109 & 21 & 2 & 5.5 \\
\texttt{pc4} & 1049 & 1458 & 37 & 2 & 7.2 \\
\texttt{ozone-level-8hr} & 1487 & 2534 & 72 & 2 & 14.8 \\
\texttt{sick} & 38 & 3772 & 29 & 2 & 15.3 \\
\texttt{abalone} & 183 & 4177 & 8 & 28 & 689.0 \\
\texttt{yeast} & 181 & 1484 & 8 & 10 & 92.6 \\
\texttt{oil\_spill} & 1018 & 8844 & 56 & 2 & 14.6 \\
\texttt{mammography} & 310 & 11183 & 6 & 2 & 42.0 \\
\texttt{wine\_quality} & 287 & 6497 & 11 & 7 & 567.2 \\
\texttt{thyroid} & 40474 & 2800 & 26 & 5 & 52.6 \\
\texttt{satimage} & 182 & 6430 & 36 & 6 & 2.4 \\
\texttt{creditcard} & 1597 & 50000 & 29 & 2 & 601.4 \\
\texttt{webpage} & 350 & 34780 & 300 & -- & -- \\
\texttt{solar-flare} & 40687 & 1066 & 12 & 6 & 7.7 \\
\texttt{iris} & sklearn & 150 & 4 & 3 & 1.0 \\
\texttt{wine} & sklearn & 178 & 13 & 3 & 1.5 \\
\texttt{breast\_cancer} & sklearn & 569 & 30 & 2 & 1.7 \\
\texttt{digits} & sklearn & 1797 & 64 & 10 & 1.1 \\
\texttt{glass} & 41 & 214 & 9 & 6 & 8.4 \\
\texttt{vowel} & 307 & 990 & 12 & 11 & 1.0 \\
\texttt{letter} & 6 & 20000 & 16 & 26 & 1.1 \\
\texttt{optdigits} & 28 & 5620 & 64 & 10 & 1.0 \\
\texttt{page-blocks} & 30 & 5473 & 10 & 5 & 175.5 \\
\texttt{sonar} & 40 & 208 & 60 & 2 & 1.1 \\
\texttt{eeg-eye-state} & 1471 & 14980 & 14 & 2 & 1.2 \\
\texttt{waveform} & 60 & 5000 & 40 & 3 & 1.0 \\
\texttt{madelon} & 1485 & 2600 & 500 & 2 & 1.0 \\
\texttt{GesturePhaseSegmentationProcessed} & 4538 & 9873 & 32 & 5 & 3.0 \\
\texttt{JapaneseVowels} & 375 & 9961 & 14 & 9 & 2.1 \\
\texttt{heart-statlog} & 53 & 270 & 13 & 2 & 1.2 \\
\texttt{ionosphere} & 59 & 351 & 34 & 2 & 1.8 \\
\texttt{wine-quality-white} & 40498 & 4898 & 11 & 7 & 439.6 \\
\texttt{wine-quality-red} & 40691 & 1599 & 11 & 6 & 68.1 \\
\texttt{mfeat-factors} & 12 & 2000 & 216 & 10 & 1.0 \\
\texttt{mfeat-fourier} & 14 & 2000 & 76 & 10 & 1.0 \\
\texttt{mfeat-karhunen} & 16 & 2000 & 64 & 10 & 1.0 \\
\texttt{mfeat-morphological} & 18 & 2000 & 6 & 10 & 1.0 \\
\texttt{mfeat-pixel} & 20 & 2000 & 240 & 10 & 1.0 \\
\texttt{mfeat-zernike} & 22 & 2000 & 47 & 10 & 1.0 \\
\texttt{semeion} & 1501 & 1593 & 256 & 10 & 1.0 \\
\texttt{synthetic\_control} & 377 & 600 & 60 & 6 & 1.0 \\
\texttt{adult} & 179 & 48842 & 14 & 2 & 3.2 \\
\texttt{splice} & 46 & 3190 & 60 & 3 & 2.2 \\
\texttt{USPS} & 41082 & 9298 & 256 & 10 & 2.2 \\
\end{longtable}

